\chapter{Fixings and Flows}\label{ch:m2:fixings-and-flows}

Between December 2007 and January 2013, euro--dollar traders at four of the
world's largest banks, members of a chat room they called ``The Cartel'', used
coded language to coordinate their trading around two moments of the day: 1:15
in the afternoon in London, when the European Central Bank's reference rate was
taken, and 4:00, when WM/Reuters struck the benchmark rate used to price the
orders of many large customers. They
told each other their clients' orders, traded together into the fix, and held
back bids or offers to protect each other's positions. In May 2015 the four
banks agreed to plead guilty in the United States and to pay more than USD~2.5
billion in criminal fines, after regulators in London and Washington had fined
five banks more than USD~3 billion between them in November 2014. This chapter
explains what a \omterm{def:m2:fixings-and-flows:fix}{benchmark fix} is, why so much currency changes hands at it, the
month-end and option-expiry flows that meet there, and what was done wrong and
what was changed.

\section{The benchmark fix}

\begin{definition}[Benchmark fix, fixing window, fixing order]\label{def:m2:fixings-and-flows:fix}
A \emph{benchmark fix}\index{benchmark fix} is an exchange rate published at a
set time of day by an administrator, from trades and quotes observed in a short
period around that time, the \emph{fixing window}\index{fixing window}. A
\emph{fixing order}\index{fixing order} is a client's order to buy or sell an
amount of currency at the fix, whatever it turns out to be; the dealer that
accepts it guarantees the rate and bears the cost of obtaining it.
\end{definition}

The most used is the WM/Reuters 4pm London closing rate
(\cref{fig:m2:fixings-and-flows:clock}), which, in the US
regulators' words, is used to price orders for many large customers and to price
derivatives around the world. A manager whose portfolio is valued against a
benchmark at that rate (One Quant Book 1, chapter 3) has a reason to trade
exactly at it: a trade done at the fix adds no tracking error. Until
February 2015 the fix was computed from trades in a one-minute window around
16:00; following recommendations of the Financial Stability Board in September
2014, WM widened it on 15 February 2015 to five minutes, from 2.5 minutes before
the hour to 2.5 minutes after. The ECB publishes its own euro reference rates
around 16:00 Central European time from a procedure around 14:10, and says that
using them for transactions is ``strongly discouraged''.

\begin{omfigure}
\centering
\begin{tikzpicture}[font=\small,x=0.62cm]
\draw[thick,-{Stealth}] (0,0) -- (19,0) node[right,font=\footnotesize] {London time};
\foreach \h in {0,2,...,18} {\draw (\h,0.1) -- (\h,-0.1) node[below,font=\scriptsize] {\h:00};}
\draw[omDef,thick] (6,0) -- (6,0.8);
\node[omDef,font=\footnotesize,anchor=south,align=center] at (6,0.8) {Tokyo option cut\\(15:00 Tokyo)};
\draw[omDef,thick] (15,0) -- (15,-0.75);
\node[omDef,font=\footnotesize,anchor=north,align=center] at (15,-0.75) {New York option cut\\(10:00 New York)};
\draw[omProp,thick] (13.17,0) -- (13.17,2.25);
\node[omProp,font=\footnotesize,anchor=south east,align=right] at (13.17,2.0) {ECB reference rate procedure\\(14:10 CET)};
\draw[omThm,thick] (16,0) -- (16,2.25);
\node[omThm,font=\footnotesize,anchor=south west,align=left] at (16.15,2.25) {WM/Reuters\\4pm fix};
\fill[omThm!25] (15.958,0) rectangle (16.042,0.35);
\end{tikzpicture}

\omcaption{The day's reference moments in London time, in winter. The WM/Reuters
fix is taken over five minutes around 16:00; the ECB's reference rates come from
a procedure around 14:10 Central European time; FX options expire most often at
10:00 New York or 15:00 Tokyo. The clocks of summer time shift some of them by an
hour. Schematic.}\label{fig:m2:fixings-and-flows:clock}
\end{omfigure}

The dealer's position is delicate. A client asks it to buy EUR~1 billion at the
fix. The dealer does not know the fix in advance; it must buy the euros itself,
and its own buying pushes the price up. If it buys them evenly during the
window, it pays on average about the fix and makes nothing but its fee. If it
buys some of them before the window, it pays the lower prices before its own
buying moved the market, then its remaining purchases and everyone else's push
the fix up, and it sells to the client at the higher fix.

\begin{proposition}[What pre-hedging a fixing order earns]\label{prop:m2:fixings-and-flows:prehedge}
Let the dealer buy $Q$ with linear permanent impact $\lambda$ per unit, a share
$p$ before the window and the rest evenly during it, and let the fix be the
window's average price. Then, from a starting price $P_0$ and with no other
price moves,
\[
\text{fix} = P_0 + \frac{\lambda Q\,(1+p)}{2}, \qquad
\text{average cost} = P_0 + \frac{\lambda Q}{2}, \qquad
\text{profit} = \frac{\lambda Q^2\,p}{2}.
\]
The profit is paid by the client, whose fix is higher by $\lambda Q p/2$ than
with no pre-hedging.
\end{proposition}

\begin{proof}
Before the window the dealer buys $pQ$ at prices rising from $P_0$ to
$P_0 + \lambda pQ$, on average $P_0 + \lambda pQ/2$. In the window it buys $(1-p)Q$
at prices rising from $P_0 + \lambda pQ$ to $P_0 + \lambda Q$, whose average is the
fix, $P_0 + \lambda Q(1+p)/2$. The average cost is $p(P_0 + \lambda pQ/2) +
(1-p)\bigl(P_0 + \lambda Q(1+p)/2\bigr) = P_0 + \lambda Q/2$, whatever $p$. The
profit is $Q$ times the difference.
\end{proof}

\begin{omfigure}
\begin{tikzpicture}
\begin{axis}[width=11cm,height=5cm,
  xlabel={minutes from 16:00},ylabel={EURUSD, pips above start},
  tick label style={font=\small},label style={font=\small},
  xmin=-10,xmax=5,ymin=-0.2,ymax=3.6,grid=major,grid style={gray!25},
  legend columns=3,legend style={font=\footnotesize,at={(0.5,-0.3)},anchor=north,draw=none,fill=none,/tikz/every even column/.append style={column sep=8pt}},
  table/col sep=comma]
\fill[omThm!10] (axis cs:-2.5,-0.2) rectangle (axis cs:2.5,3.6);
\node[font=\scriptsize,omThm] at (axis cs:0,3.35) {fixing window};
\addplot[omDef,very thick] table[x=t,y=p0]{figdata/markets-2/17-fixings-and-flows/path.csv};
\addplot[omProp,very thick,dashed] table[x=t,y=p50]{figdata/markets-2/17-fixings-and-flows/path.csv};
\addplot[omThm,very thick,dotted] table[x=t,y=p100]{figdata/markets-2/17-fixings-and-flows/path.csv};
\draw[omDef,thin] (axis cs:-2.5,1.5) -- (axis cs:2.5,1.5);
\draw[omThm,thin] (axis cs:-2.5,3.0) -- (axis cs:2.5,3.0);
\legend{all in the window (fix $+1.5$),half pre-hedged (fix $+2.25$),all pre-hedged (fix $+3$)}
\end{axis}
\end{tikzpicture}

\omcaption{The price of EURUSD while a dealer buys EUR~1 billion for a client's
\omterm{def:m2:fixings-and-flows:fix}{fixing order}, with a permanent impact of 0.3 pips per EUR~100 million, when it
buys all of it in the five-minute window, half before, or all before. The thin
lines mark the fix in the first and last cases. The dealer's average cost is 1.5
pips in each case; the fix is not. Illustrative; data: the chapter's
tutorial.}\label{fig:m2:fixings-and-flows:path}
\end{omfigure}

\begin{example}[A billion euros at the fix]\label{ex:m2:fixings-and-flows:order}
With $\lambda = 0.3$ pips per EUR~100 million, buying EUR~1 billion moves the price
3 pips. Bought evenly in the window, the fix is 1.5 pips above the start and the
dealer's average cost is the same. Half bought before the window, the fix is 2.25
pips up and the dealer earns 0.75 pips on EUR~1 billion, USD~75\,000; all before,
USD~150\,000. The client pays the same amounts more than with no pre-hedging.
\end{example}

\section{Month-end and index rebalancing flows}

\begin{definition}[Month-end rebalancing]\label{def:m2:fixings-and-flows:monthend}
\emph{Month-end rebalancing}\index{month-end rebalancing} is the trading that
investors do at the end of each month, some of it at the 4pm fix, to restore
currency hedges and portfolio weights that the month's price moves have
disturbed.
\end{definition}

Two mechanisms make month-end flows predictable. A foreign investor that hedges
half of its dollar equities finds, after a month in which they rose 4\%, that its
hedge covers less than half: it sells more dollars forward, half of 4\% of its
holdings. A fund that keeps 60\% in American assets and 40\% in European ones
finds, after a month in which the first rose 5\% and the second fell 1\%, that it
holds too much of the first; it sells American assets, and dollars, and buys
European ones. Both flows are known in sign, and roughly in size, days in
advance (\cref{fig:m2:fixings-and-flows:hedgeflow}).

\begin{omfigure}
\begin{tikzpicture}
\begin{axis}[width=11cm,height=4.8cm,
  xlabel={US equity return over the month (\%)},ylabel={dollar sales (USD bn)},
  tick label style={font=\small},label style={font=\small},
  xmin=-6,xmax=6,ymin=-95,ymax=95,grid=major,grid style={gray!25},
  legend columns=3,legend style={font=\footnotesize,at={(0.5,-0.3)},anchor=north,draw=none,fill=none,/tikz/every even column/.append style={column sep=8pt}},
  table/col sep=comma]
\addplot[omDef,thick,dashed] table[x=ret,y=h25]{figdata/markets-2/17-fixings-and-flows/hedgeflow.csv};
\addplot[omThm,very thick] table[x=ret,y=h50]{figdata/markets-2/17-fixings-and-flows/hedgeflow.csv};
\addplot[omProp,thick,dotted] table[x=ret,y=h75]{figdata/markets-2/17-fixings-and-flows/hedgeflow.csv};
\legend{hedged 25\%,hedged 50\%,hedged 75\%}
\end{axis}
\end{tikzpicture}

\omcaption{Month-end dollar sales (negative: purchases) needed to keep the hedge
of USD~2 trillion of US equities at its ratio, by the month's equity return and
the hedge ratio. The flow is known in sign as soon as the month's return is.
Illustrative; data: the chapter's tutorial.}\label{fig:m2:fixings-and-flows:hedgeflow}
\end{omfigure}

\begin{example}[Two month-end flows]\label{ex:m2:fixings-and-flows:flows}
Investors holding USD~2 trillion of American equities, hedged at 50\%, sell
$0.5 \times 4\% \times 2$ trillion $=$ USD~40 billion forward after a 4\% rise. A
fund of USD~100 billion at 60/40, after the American part gained 5\% and the
European part lost 1\%, must sell USD~1.44 billion of American assets and buy the
same value of European ones.
\end{example}

Anyone can make these estimates, so flows that are known in advance may be
traded ahead: a dealer or hedge fund that expects dollar sales at the fix can
sell first, and the price may move before the fix rather than at it. What the
investor pays for the certainty of a benchmark rate is therefore uncertain, a
reason to spread large month-end trades over a longer period or across
execution algorithms.

\begin{dated}{2026-09}{The size of the stock}\label{dat:m2:fixings-and-flows:size}
Holdings of US corporate equities by the rest of the world were valued at
USD~22.2 trillion in the second quarter of 2026 (Federal Reserve Financial
Accounts). No official figure gives the share of it that is currency-hedged,
and hedge ratios differ by investor and country.
\end{dated}

\section{Option expiries and the cuts}

\begin{definition}[Option cut]\label{def:m2:fixings-and-flows:cut}
An \emph{option cut}\index{option cut} is the time of day at which an FX option
expires: most commonly 10:00 New York time (the New York cut) or 15:00 Tokyo
time (the Tokyo cut).
\end{definition}

Large options expiring at a cut with strikes near spot create flows of their own.
A dealer long an option hedges it by selling as the price rises and buying as it
falls, which pins the price near the strike as expiry approaches; a dealer short
it does the opposite and adds to moves. At the cut the option's delta jumps to
zero or one and its hedge must be unwound or completed. Traders watch lists of
large expiries near spot for the day's cut; the options themselves are the subject
of \cref{ch:m2:the-fx-options-market}.

\section{The fixing scandal and what changed}

\begin{definition}[Banging the close]\label{def:m2:fixings-and-flows:banging}
\emph{Banging the close}\index{banging the close} is concentrating trades in the
moments that set a benchmark or settlement price, in order to move it in a
direction that benefits a position.
\end{definition}

The regulators' findings describe the mechanism of
\cref{prop:m2:fixings-and-flows:prehedge} turned into collusion. In November
2014 the UK Financial Conduct Authority fined Citibank, HSBC, JPMorgan Chase, the
Royal Bank of Scotland and UBS GBP~1.11 billion for failing to control their G10
spot FX businesses between January 2008 and October 2013, during which traders
shared confidential information about clients' orders and attempted to
manipulate rates, including in collusion with traders at other firms; the US
Commodity Futures Trading Commission fined the same five banks more than USD~1.4
billion for attempted manipulation of FX benchmarks, primarily the 4pm fix. In
May 2015 Citicorp, JPMorgan, Barclays and the Royal Bank of Scotland agreed to
plead guilty to a US antitrust conspiracy to fix prices and rig bids in euros and
dollars, with fines from USD~395 million to USD~925 million each.

\begin{omfigure}
\begin{tikzpicture}
\begin{axis}[width=11cm,height=4.8cm,
  xlabel={share of the order bought before the window, $p$},ylabel={dealer profit (USD thousand)},
  tick label style={font=\small},label style={font=\small},
  xmin=0,xmax=1,ymin=-80,ymax=380,grid=major,grid style={gray!25},
  legend columns=2,legend style={font=\footnotesize,at={(0.5,-0.3)},anchor=north,draw=none,fill=none,/tikz/every even column/.append style={column sep=8pt}},
  table/col sep=comma]
\addplot[name path=hi,draw=none,forget plot] table[x=p,y=hi]{figdata/markets-2/17-fixings-and-flows/pnl.csv};
\addplot[name path=lo,draw=none,forget plot] table[x=p,y=lo]{figdata/markets-2/17-fixings-and-flows/pnl.csv};
\addplot[omThm!20,forget plot] fill between[of=hi and lo];
\addplot[omThm,very thick,mark=*,mark size=1.2pt] table[x=p,y=pnl]{figdata/markets-2/17-fixings-and-flows/pnl.csv};
\legend{expected profit (band: $\pm$ one standard deviation)}
\end{axis}
\end{tikzpicture}

\omcaption{The dealer's profit on a EUR~1 billion \omterm{def:m2:fixings-and-flows:fix}{fixing order} by the share
bought before the window, with the chapter's impact and a random move of two
pips (standard deviation) between the pre-hedge and the window. The expected
profit, taken from the client, grows with pre-hedging, and so does its risk.
Illustrative; data: the chapter's tutorial.}\label{fig:m2:fixings-and-flows:pnl}
\end{omfigure}

The fix changed with them. The wider window makes the average harder to move,
and WM added price sources. The \omterm{def:m2:fx-spot-microstructure:code}{FX Global Code}
(\cref{ch:m2:fx-spot-microstructure}) says that a dealer handling \omterm{def:m2:fixings-and-flows:fix}{fixing orders}
should not share information inappropriately or try to influence the rate,
whether by collusion or otherwise, should not intentionally influence the fix
to benefit from it, and should price such orders transparently and consistently
with the risk it bears. Its Principle~11 allows pre-hedging only by a dealer
acting as principal, to manage the risk of an anticipated client order in a way
designed to benefit the client, and requires the practice to be explained to
clients. Trading ahead of a \omterm{def:m2:fixings-and-flows:fix}{fixing order} is therefore neither banned nor free:
whether it serves the client is the test, and the line is argued over still.

\section{Tutorial: month-end flows and a fixing order}\label{tut:m2:fixings-and-flows:flows}

\begin{tutorial}
\textbf{Goal.} Estimate month-end hedge and portfolio rebalancing flows, and
price a \omterm{def:m2:fixings-and-flows:fix}{fixing order} with and without pre-hedging. \textbf{End state:}
\cref{fig:m2:fixings-and-flows:path,fig:m2:fixings-and-flows:pnl},
\cref{ex:m2:fixings-and-flows:order,ex:m2:fixings-and-flows:flows} and the
numbers of the weekend problem.
\begin{tutsteps}
\item \textbf{The flows}: hedge and weight rebalancing.
\omcode{code/firm/fixflow/firm_fixflow.py}{14}{23}{Month-end hedge and portfolio rebalancing flows.}\label{lst:m2:fixings-and-flows:flows}
\item \textbf{The \omterm{def:m2:fixings-and-flows:fix}{fixing order}}: the closed form of
\cref{prop:m2:fixings-and-flows:prehedge}, and a simulation with a random move.
\omcode{code/firm/fixflow/firm_fixflow.py}{26}{51}{Fix, dealer cost and profit of a fixing order.}\label{lst:m2:fixings-and-flows:fix}
\item \textbf{Run} \texttt{fixflow\_demo.month\_end()},
\texttt{fixflow\_demo.fixing\_order} for $p = 0$, 0.5 and 1, and
\texttt{fig\_fixflow.py}.
\end{tutsteps}
\textbf{What to change next.} Add temporary impact, which decays after each
trade, and show that it makes buying inside the window cost more than the fix;
then find the pre-hedge share that minimises the dealer's risk for a fixed
expected profit of zero.
\end{tutorial}

\section{Build: the month-end flow estimator}\label{bld:m2:fixings-and-flows:fixflow}

\begin{build}
\bfield{Purpose} The miniature firm handles clients' \omterm{def:m2:fixings-and-flows:fix}{fixing orders} and trades
around month ends: it must estimate the flows it will face and know, and
disclose, the economics of how it executes \omterm{def:m2:fixings-and-flows:fix}{fixing orders}.
\bfield{Interface} \texttt{hedge\_rebalance(value\_usd, asset\_return,
hedge\_ratio)}; \texttt{weight\_rebalance(values, targets)};
\texttt{fix\_and\_cost(q, lam, p, p0)}; \texttt{dealer\_pnl(q, lam, p)};
\texttt{simulate\_pnl(q, lam, p, sigma\_pre, n, seed)}.
\bfield{Rules} Dollar sales positive; linear permanent impact; the fix as the
window's average price; the random move applied after the pre-hedge.
\bfield{Acceptance tests} \texttt{code/firm/fixflow/tests/}: signs of the hedge
flow; the fund sells its outperforming region; with no pre-hedge the dealer buys
at the fix; the cost independent of $p$; pre-hedging adds risk.
\bfield{Stretch} Flows by country from public holdings data and monthly index
returns; temporary impact; the fix computed from a simulated order book as WM
computes it, as a median of observed rates.
\end{build}

\begin{omsources}
\item US Department of Justice, ``Five major banks agree to parent-level guilty
pleas'', 20 May 2015.
\item Financial Conduct Authority, press release, 12 November 2014; Commodity
Futures Trading Commission, press release 7056-14, 12 November 2014.
\item Financial Stability Board, \emph{Foreign Exchange Benchmarks: Final
Report}, September 2014; WM/Reuters presentation to the ECB FX Contact Group,
November 2015.
\item European Central Bank, euro foreign exchange reference rates; Federal
Reserve, Financial Accounts of the United States.
\end{omsources}

\section{Exercises}

\begin{exercise}[$\star$]\label{exo:m2:fixings-and-flows:1}
Why would an index fund want to trade currency at the 4pm fix rather than earlier
in the day at a better price?
\end{exercise}

\begin{exercise}[$\star$]\label{exo:m2:fixings-and-flows:2}
A Japanese investor holds USD~50 billion of US equities hedged at 70\%. They rise
3\% over the month. What does it trade at month-end?
\end{exercise}

\begin{exercise}[$\star$]\label{exo:m2:fixings-and-flows:3}
What changed in the WM/Reuters fix in February 2015, and why does it make
manipulation harder?
\end{exercise}

\begin{exercise}[$\star\star$]\label{exo:m2:fixings-and-flows:4}
In \cref{ex:m2:fixings-and-flows:order}, give the fix and the dealer's profit if
it pre-hedges 30\%.
\end{exercise}

\begin{exercise}[$\star\star$]\label{exo:m2:fixings-and-flows:5}
A fund at 60/40 US/Europe with USD~100 billion sees the American part fall 3\%
and the European part rise 2\% in dollars. What does it trade?
\end{exercise}

\begin{exercise}[$\star\star$]\label{exo:m2:fixings-and-flows:6}
A dealer is long a large EURUSD option struck just above spot, expiring at the
New York cut. How does its hedging affect the price in the morning?
\end{exercise}

\begin{exercise}[$\star\star\star$]\label{exo:m2:fixings-and-flows:7}
\emph{Coding.} With \texttt{simulate\_pnl}, give the mean and standard deviation of
the dealer's profit at $p = 0.5$ and $p = 1$ with the chapter's parameters.
\end{exercise}

\begin{exercise}[$\star\star\star$]\label{exo:m2:fixings-and-flows:8}
\emph{Find the flaw.} ``The dealer's buying before the window is just hedging:
its average cost is the same whatever it does, so the client is not harmed.''
Correct it.
\end{exercise}

\section{Problem: Guaranteed at the Fix}

\begin{problem}[{Weekend problem --- what a fixing order earns a dealer}]\label{pb:m2:fixings-and-flows:1}
A pension fund asks a dealer to buy EUR~1 billion at the 4pm fix at month end.
The dealer estimates a permanent impact of 0.3 pips per EUR~100 million, and a
random move of two pips (standard deviation) between 15:50 and the window.

\medskip\noindent\textbf{Part I --- No pre-hedging.}
\begin{enumerate}
\item How far does the dealer's buying move EURUSD?
\item Give the fix, relative to the start.
\item Give the dealer's average cost and profit.
\item What risk does the dealer bear?
\item What does the dealer charge for the service, and how?
\end{enumerate}

\medskip\noindent\textbf{Part II --- Pre-hedging.}
\begin{enumerate}[resume]
\item The dealer buys half before 15:57:30. Give the fix.
\item Give its expected profit, and who pays it.
\item Give the standard deviation of its profit.
\item Give the same for full pre-hedging.
\item Why does the average cost not depend on the share pre-hedged?
\end{enumerate}

\medskip\noindent\textbf{Part III --- Other flows.}
\begin{enumerate}[resume]
\item Other clients sell EUR~600 million at the same fix. What changes?
\item Month-end hedge rebalancing adds USD~40 billion of dollar sales. Which way does
it push EURUSD, and when?
\item Why might the price move before the window rather than in it?
\item How would the fund reduce its cost?
\item What does the \omterm{def:m2:fx-spot-microstructure:code}{FX Global Code} require of the dealer?
\end{enumerate}

\medskip\noindent\textbf{Part IV --- Judgement.}
\begin{enumerate}[resume]
\item When is pre-hedging legitimate risk management?
\item What made the conduct of 2008--2013 criminal?
\item Why does a five-minute window help, and what does it not prevent?
\item State the \emph{named result}: the dealer's expected profit on the order with
and without pre-hedging.
\item In one sentence: who pays when a dealer trades ahead of a \omterm{def:m2:fixings-and-flows:fix}{fixing order}?
\end{enumerate}
\end{problem}

\section{Interview questions}

\begin{interviewq}[$\star$ \iqroles{trader, bank}]\label{iq:m2:fixings-and-flows:1}
What is the WM/Reuters fix, and who uses it?
\end{interviewq}

\begin{interviewq}[$\star$ \iqroles{trader, researcher}]\label{iq:m2:fixings-and-flows:2}
Why are month-end FX flows predictable, and in which direction are they after
a strong month for US equities?
\end{interviewq}

\begin{interviewq}[$\star\star$ \iqroles{trader}]\label{iq:m2:fixings-and-flows:3}
A client gives you a EUR~2 billion \omterm{def:m2:fixings-and-flows:fix}{fixing order}. How do you execute it, and what
do you tell the client?
\end{interviewq}

\begin{interviewq}[$\star\star$ \iqroles{researcher}]\label{iq:m2:fixings-and-flows:4}
Show that pre-hedging a \omterm{def:m2:fixings-and-flows:fix}{fixing order} earns the dealer money in a linear-impact
model. Where does the money come from?
\end{interviewq}

\begin{interviewq}[$\star\star$ \iqroles{bank, trader}]\label{iq:m2:fixings-and-flows:5}
What were the FX benchmark cases of 2014 and 2015 about, and what changed as a
result?
\end{interviewq}

\begin{interviewq}[$\star\star\star$ \iqroles{developer, researcher}]\label{iq:m2:fixings-and-flows:6}
Design a surveillance system that detects \omterm{def:m2:fixings-and-flows:banging}{banging the close} around the 4pm fix.
\end{interviewq}
